\unsection{ВИСНОВКИ}

У курсовій роботі розроблено контейнеризовану пошукову систему, що демонструє паралелізм на
всіх рівнях: паралельну побудову інвертованого індексу власним пулом потоків, потокобезпечну
спільну структуру індексу у двох варіантах -- неблокуючому та блокувальному, -- безблокувальне
читання незмінного знімка індексу і паралельну обробку запитів багатьох клієнтів власним
HTTP/1.1- і WebSocket-сервером. Пошуковий конвеєр розширює запит через локальну модель
qwen2.5:1.5b і ранжує документи функцією BM25. Система запускається командою
\code{docker compose up}, повністю реалізує фіксований контракт REST API і WebSocket, має два
клієнти (CLI і Web UI) і написана мовою Java 21 без жодних сторонніх бібліотек. Її
коректність підтверджують 38 тестів, зокрема стресові порівняння паралельних побудов з
послідовним еталоном і наскрізні тести під навантаженням.

Результати дослідження масштабованості на процесорі з 6 ядрами і 12 апаратними потоками:
\begin{enumerate}[label=\arabic*)]
\item Час побудови індексу корпусу з 16\,793 документів (1,62~ГБ) скорочується з 10,12~с на
одному потоці до 2,09~с на десяти: прискорення 4,85. До шести потоків (кількість фізичних
ядер) ефективність становить 0,67--0,87, після шести додаткові потоки використовують SMT і
дають лише 20~\% приросту, а на 16 потоках час трохи зростає.
\item Апроксимація законом Амдала дає послідовну частку 0,125, тоді як справді послідовні
ділянки займають менше 2~\% часу; експериментальна частка Карпа -- Флатта зростає з
кількістю потоків. Отже, масштабування обмежують накладні витрати, що зростають із
паралелізмом: сумарний час токенізації, яка не має жодної синхронізації, на 12 потоках у
2,25 раза більший, ніж на одному. Причина -- конкуренція за спільні апаратні ресурси (ядра
при SMT, кеш L3, шина пам'яті, частота процесора).
\item Заміна токенізації, що створювала рядок на кожне слово, байтовою хеш-таблицею задачі
прискорила побудову в 1,28--1,48 раза за будь-якої кількості потоків. Зменшення виділення
пам'яті в гарячому циклі дало більший ефект, ніж будь-яка зміна механізму синхронізації.
\item З однією партицією обидві реалізації спільної структури перестають масштабуватися після
4 потоків (3,4--3,7~с проти 2,1~с). Блокувальна при цьому сумарно простоює на замку 22,4~с,
неблокуюча не простоює зовсім, але виконує злиття тим самим послідовним потоком-власником:
неблокуючий алгоритм усуває очікування, а не серіалізацію. Від 8 партицій синхронізація
перестає впливати на час, і обидві реалізації показують однаковий результат.
\item Розміщення корпусу в томі Docker замість каталогу диска Windows скоротило побудову на
12 потоках з 42,3 до 4,0~с: через міст файлових систем Windows -- WSL2 задача була обмежена
введенням-виведенням і не могла показати масштабування.
\item Пропускна здатність пошуку зростає майже лінійно до 8 клієнтів і досягає 12\,004
запитів/с на 16 клієнтах за медіанної латентності менше 1~мс; далі вона обмежена кількістю
апаратних потоків, а латентність зростає. Пошукові потоки не синхронізуються між собою, бо
читають незмінний знімок індексу.
\item Розширення запиту через LLM на CPU обмежує пропускну здатність рівнем 2,3--2,4 запиту/с
незалежно від кількості клієнтів: модель є послідовним ресурсом, і латентність зростає за
законом Літтла. Кеш розширень повертає пропускну здатність повторних запитів до
8\,634 запитів/с, а дедлайн і підміна знімка після розширення не дають очікуванню на
модель утримувати ресурси сервера.
\item Готовий індекс займає близько 350~МіБ купи, витоку пам'яті між перебудовами немає, але
під час перебудови за схемою build-then-swap потрібна пам'ять для двох індексів і проміжних
таблиць, тож пікове використання становить кілька гігабайтів.
\end{enumerate}

Подальші напрями покращення: зменшити обсяг пам'яті, що виділяється під час побудови
(повторно використовувати таблиці задач, стискати списки входжень дельта-кодуванням), що
зменшить конкуренцію за шину пам'яті; не розширювати запит, якщо жоден його терм не знайдено
в індексі; винести модель на GPU або обробляти розширення пакетами, щоб послабити найбільше
послідовне обмеження конвеєра; врахувати в ранжуванні близькість слів запиту в документі.
